% part: intuitionistic-logic
% chapter: propositions-as-types
% section: types

\documentclass[../../../include/open-logic-section]{subfiles}

\begin{document}

\olfileid{int}{pty}{typ}

\olsection{في الأنماط}

لا يعيّن حد برهان مثل~$(MN)$ بمفرده البرهان الذي يرمّزه تعيينًا وحيدًا؛
ذلك أن المتغيرات الحرة، وهي حدود برهان البديهيات، لا تحدد~!!{formula}s في
البديهيات المقابلة. أما إذا حددناها في سياق~$\OLId{greek-variable}{Gamma}$، كان البرهان، متى
وُجد، \emph{معيّنًا تعيينًا وحيدًا} بالنسبة إلى ذلك السياق.

\begin{defn}
\emph{سياق} حد برهان~$\OLId{latin-looped}{N}$ هو مجموعة~$\OLId{greek-variable}{Gamma}$ تحتوي، لكل متغير حر~$\OLId{latin-ordinary}{x}$ في~$\OLId{latin-looped}{N}$،
زوجًا واحدًا بالضبط~$\OLId{latin-ordinary}{x}:!A$.
\end{defn}

وتنبّه إلى أن التعريف لا يشترط أن تحتوي~$\OLId{greek-variable}{Gamma}$ أزواجًا~$\OLId{latin-ordinary}{x}:!A$
\emph{للمتغيرات الحرة في~$\OLId{latin-looped}{N}$ وحدها}. ولذلك، مثلًا، تكون
$\OLId{greek-variable}{Gamma}=\{\OLId{latin-ordinary}{x}:!A,\OLId{latin-ordinary}{y}:!B,\OLId{latin-ordinary}{z}:!C\}$ سياقًا لـ~$\pair{\OLId{latin-ordinary}{x}}{\OLId{latin-ordinary}{y}}$.

\begin{defn}\ollabel{defn:type}
  لنفرض أن~$\OLId{greek-variable}{Gamma}$ سياق لـ~$\OLId{latin-looped}{N}$. يُعرّف \emph{نوع}~$\OLId{latin-looped}{N}$ بالنسبة إلى
  $\OLId{greek-variable}{Gamma}$ استقرائيًا كما يأتي:
  \begin{enumerate}
  \item نوع~$\OLId{latin-ordinary}{x}$ هو~$!A$ إذا وفقط إذا كان~$\OLId{latin-ordinary}{x}:!A\in\OLId{greek-variable}{Gamma}$.
  \item إذا كان نوع~$\OLId{latin-looped}{N}$ بالنسبة إلى~$\OLId{latin-ordinary}{x}:!A,\OLId{greek-variable}{Gamma}$ هو~$!B$، فنوع
    $\lambd[\typeof{\OLId{latin-ordinary}{x}}{!A}][\OLId{latin-looped}{N}]$ هو~$!A\lif !B$.
  \item إذا كان نوع~$\OLId{latin-looped}{N}$ هو~$!A\lif !B$ ونوع~$\OLId{latin-looped}{M}$ هو~$!A$، فنوع~$(NM)$
    هو~$!B$.
  \item إذا كان نوع~$\OLId{latin-looped}{N}$ هو~$!A$ ونوع~$\OLId{latin-looped}{M}$ هو~$!B$، فنوع~$\pair{\OLId{latin-looped}{N}}{\OLId{latin-looped}{M}}$
    هو~$!A\land !B$.
  \item إذا كان نوع~$\OLId{latin-looped}{N}$ هو~$!A_\OLMathNumeral{1}\land !A_\OLMathNumeral{2}$، فنوع~$\proj{\OLId{latin-ordinary}{i}}{\OLId{latin-looped}{N}}$ هو
    $!A_\OLId{latin-ordinary}{i}$.
  \item إذا كان نوع~$\OLId{latin-looped}{N}$ هو~$!A$، فنوع~$\inj[!B]{\OLId{latin-ordinary}{i}}{\OLId{latin-looped}{N}}$ هو
    $!A\lor !B$ إذا كان~$\OLId{latin-ordinary}{i}=\OLMathNumeral{1}$، وهو~$!B\lor !A$ إذا كان~$\OLId{latin-ordinary}{i}=\OLMathNumeral{2}$.
  \item إذا كان نوع~$\OLId{latin-looped}{M}$ هو~$!A\lor !B$، وكان نوع~$\OLId{latin-looped}{N}_\OLMathNumeral{1}$ بالنسبة إلى
    $\OLId{latin-ordinary}{x}_\OLMathNumeral{1}:!A,\OLId{greek-variable}{Gamma}$ هو~$!C$، ونوع~$\OLId{latin-looped}{N}_\OLMathNumeral{2}$ بالنسبة إلى~$\OLId{latin-ordinary}{x}_\OLMathNumeral{2}:!B,\OLId{greek-variable}{Gamma}$
    هو~$!C$، فنوع~$\dcase{\OLId{latin-looped}{M}}{\OLId{latin-ordinary}{x}_\OLMathNumeral{1}}{\OLId{latin-looped}{N}_\OLMathNumeral{1}}{\OLId{latin-ordinary}{x}_\OLMathNumeral{2}}{\OLId{latin-looped}{N}_\OLMathNumeral{2}}$ هو~$!C$.
  \item إذا كان نوع~$\OLId{latin-looped}{N}$ هو~$\lfalse$، فنوع~$\abort{!A}{\OLId{latin-looped}{N}}$ هو~$!A$.
  \end{enumerate}
\end{defn}

\begin{prop}
لكل حد برهان~$\OLId{latin-looped}{N}$ نوع واحد على الأكثر بالنسبة إلى سياق~$\OLId{greek-variable}{Gamma}$ له؛ أي إن
نوعه، إذا وُجد، وحيد بالنسبة إلى~$\OLId{greek-variable}{Gamma}$.
\end{prop}

\begin{proof}
  بالاستقراء في~$\OLId{latin-looped}{N}$.
\end{proof}

نكتب~$\OLId{greek-variable}{Gamma}\Sequent \OLId{latin-looped}{N}:!A$ إذا كان نوع~$\OLId{latin-looped}{N}$ بالنسبة إلى~$\OLId{greek-variable}{Gamma}$ هو
$!A$. وينبغي أن تبدو بنود~\olref{defn:type} مألوفة: فهي بالضبط قواعد
\Log{tN2i} مع فرق طفيف واحد، هو أن السياق~$\OLId{greek-variable}{Gamma}$ يبقى نفسه طوال البرهان.
ويمكننا ترميز التعريف في حساب~\Log{tN3i} ذي القواعد الواردة في
\olref{tab:tN3ip}.

\subfile{rules-tN3}

حدود البرهان حدود في نسخة من حساب لامبدا تسمى \emph{حساب لامبدا المنمّط
ذو الأنماط الضربية والجمعية والخالية}. ولا يحتوي حساب لامبدا التقليدي إلا
حدودًا مبنية من المتغيرات، و\emph{تجريدات لامبدا}~$\lambd[\OLId{latin-ordinary}{x}][\OLId{latin-looped}{N}]$،
والتطبيق~$(MN)$؛ ولا يحتوي وسم~$!A$ في تجريد لامبدا المنمّط
$\lambd[\typeof{\OLId{latin-ordinary}{x}}{!A}][\OLId{latin-looped}{N}]$. وحساب لامبدا نموذج للحوسبة تام تورنغ، أي يمكن
تمثيل كل دالة جزئية قابلة للحساب فيه بحد. أما نسخته التي تعطيها حدود البرهان
وقواعد التنميط فهي نموذج حوسبة مقيّد، غير أن أفكاره الأساسية هي نفسها. ويفرض
إسناد الأنماط هذا القيد: فلا تعد حدودًا مشروعة إلا الحدود \emph{الصحيحة
النوع}، أي حدود البرهان ذات الأنماط بالنسبة إلى سياق~$\OLId{greek-variable}{Gamma}$. فمثلًا،
$\lambd[\OLId{latin-ordinary}{x}][(xx)]$ حد مشروع في حساب لامبدا التقليدي غير المنمّط، غير أن
$\lambd[\typeof{\OLId{latin-ordinary}{x}}{!A}][(xx)]$ ليس صحيح النوع لأي~$!A$. ذلك أن قواعد
التنميط تقتضي، كي يكون لـ~$(xx)$ نوع، أن يكون نوع~$\OLId{latin-ordinary}{x}$ الأول
$!A\lif !B$ ونوع~$\OLId{latin-ordinary}{x}$ الثاني~$!A$؛ وهذا محال لأن~$!A$~!!{formula} جزئية
أصلية من~$!A\lif !B$.

ويمكن فهم الأنماط حدسيًا بوصفها أصناف الأشياء التي تسميها حدود من تلك
الأنماط. فالحد من النوع~$!A\land !B$ يعيّن زوجًا مركبته الأولى من النوع
$!A$ والثانية من النوع~$!B$؛ أي~!!a{element} من الضرب~$!A\times !B$.
والحد من النوع~$!A\lif !B$ دالة من أشياء النوع~$!A$ إلى أشياء النوع~$!B$.
والحد من النوع~$!A\lor !B$ إما شيء من النوع~$!A$ وإما شيء من النوع~$!B$،
ومعه معلومة تحدد أيهما هو. وهذا يشبه الاتحاد المنفصل للمجموعتين~$!A$ و
$!B$ المعرّف بأنه
$\{\tuple{\OLId{latin-ordinary}{i},\OLId{latin-ordinary}{x}}:\OLId{latin-ordinary}{i}=\OLMathNumeral{1}\text{ أو }\OLMathNumeral{2}\text{، }\OLId{latin-ordinary}{x}\in !A\text{ إذا كان }\OLId{latin-ordinary}{i}=\OLMathNumeral{1}\text{، }
\OLId{latin-ordinary}{x}\in !B\text{ إذا كان }\OLId{latin-ordinary}{i}=\OLMathNumeral{2}\}$. والنوع~$\lfalse$ هو المجموعة الخالية؛
أي إن حدًا من هذا النوع لا يجوز أن يسمي شيئًا. ولأن الاستنباط الطبيعي، ومنه
\Log{tN3i}، متسق، فلا يوجد~!!{proof} لـ~$\lfalse$ بغير فروض مفتوحة؛ ومن
ثم لا يوجد حد برهان مغلق، أي حد بلا متغيرات حرة، نوعه~$\lfalse$.

تنتج مؤثرات حساب لامبدا المنمّط حدودًا تسمي حدسيًا الأشياء من الصنف الصحيح،
بشرط أن تسمي الحدود التي تطبق عليها أشياء معينة. فمثلًا، تقول العبارة
``إذا كان~$\OLId{latin-looped}{N}_\OLMathNumeral{1}:!A$ و$\OLId{latin-looped}{N}_\OLMathNumeral{2}:!B$، كان
$\pair{\OLId{latin-looped}{N}_\OLMathNumeral{1}}{\OLId{latin-looped}{N}_\OLMathNumeral{2}}:!A\land !B$'' إن~$\OLId{latin-looped}{N}_\OLMathNumeral{1}$ إذا سمّى شيئًا من النوع~$!A$
وسمّى~$\OLId{latin-looped}{N}_\OLMathNumeral{2}$ شيئًا من النوع~$!B$، سمّى~$\pair{\OLId{latin-looped}{N}_\OLMathNumeral{1}}{\OLId{latin-looped}{N}_\OLMathNumeral{2}}$ شيئًا من النوع
$!A\land !B$.

\end{document}
